\chapter{\texttt{detector.interaction} --- neutrino-target interaction cross sections}
\label{ch:detector-interaction}

The folding chain of \cref{ch:detector-common} treats the differential
cross section $d\sigma/dT(E_\nu,T)$ as an external input. This chapter
supplies the one interaction currently implemented: elastic neutrino-electron
scattering, the process every liquid-scintillator or water-based solar and
reactor detector (Borexino, KamLAND, Super-Kamiokande) uses to detect
sub-MeV--few-MeV neutrinos, and the only channel
\texttt{detector.borexino} currently folds (\cref{ch:detector-borexino}).

\section{\texttt{neutrino\_electron.py} --- elastic $\nu$-$e$ scattering}

\modulehead{detector/interaction/neutrino\_electron.py}{Standard Model
tree-level differential cross section for $\nu + e^- \to \nu + e^-$, no
radiative corrections.}

\begin{theorybox}
\subsection*{Two amplitudes, and why they single out the electron flavour}
Elastic neutrino-electron scattering, $\nu_\alpha + e^- \to \nu_\alpha +
e^-$, proceeds through two tree-level weak amplitudes, exactly the same two
channels that build the coherent matter potential of \cref{ch:core-common}:
\begin{itemize}
  \item \textbf{Charged current (CC), $W^\pm$ exchange.} The incoming
    neutrino and the target electron can exchange a $W$ boson only if the
    final state can still be written as (some neutrino) + (some electron);
    for a target of ordinary electrons (no muons or taus present) this is
    only possible when the incoming neutrino is itself $\nu_e$. This is the
    identical argument, applied at the level of a single scattering event
    rather than coherent forward scattering, that in \cref{ch:core-common}
    restricts the CC matter potential $V_{CC}$ to the electron entry alone.
  \item \textbf{Neutral current (NC), $Z^0$ exchange.} A $Z^0$ couples to
    all three active flavours identically (lepton universality), so this
    amplitude contributes equally for $\nu_e$, $\nu_\mu$, and $\nu_\tau$.
\end{itemize}
For $\nu_e$, the CC and NC amplitudes interfere (both have the same initial
and final state); for $\nu_\mu,\nu_\tau$ only the NC amplitude exists. The
standard way to encode this is to absorb the interference into a single pair
of \emph{effective} left/right chiral couplings $(g_L,g_R)$, so that a
single formula for $d\sigma/dT$ (below) covers both cases:
\begin{align}
  \nu_e\ \text{(CC + NC)}:\quad
    & g_L = \tfrac12 + \sin^2\theta_W, && g_R = \sin^2\theta_W, \\
  \nu_\mu,\nu_\tau\ \text{(NC only)}:\quad
    & g_L = -\tfrac12 + \sin^2\theta_W, && g_R = \sin^2\theta_W,
\end{align}
with $\sin^2\theta_W$ the (effective, low-energy) weak mixing angle
\parencite{ParticleDataGroup}. The two cases differ only by the additive
$+1$ shift of $g_L$ contributed by the CC amplitude available to $\nu_e$
alone; $g_R$, which multiplies the purely-NC right-handed electron coupling,
is identical for all three flavours. This is exactly why the
\texttt{detector.common} folding chain (\cref{ch:detector-common}) only ever
needs a two-channel split $\Phi_e/\Phi_x$: $\nu_\mu$ and $\nu_\tau$ are, by
this construction, kinematically indistinguishable in this interaction.

\subsection*{Differential cross section}
The Standard Model tree-level result for elastic scattering off an electron
initially at rest, differential in the electron recoil kinetic energy $T$,
is \parencite{BahcallKamionkowskiSirlin1995,ParticleDataGroup}
\begin{equation}
  \keyresult{\frac{d\sigma}{dT}(E_\nu,T) =
    \frac{2\,G_F^2\,m_e}{\pi}\,(\hbar c)^2\,
    \left[g_L^2 + g_R^2\left(1-\frac{T}{E_\nu}\right)^{\!2}
      - g_L g_R\,\frac{m_e\,T}{E_\nu^2}\right],}
  \label{eq:nue-dsdt}
\end{equation}
valid for $0\leq T\leq T_{\max}(E_\nu)$ and zero outside that range. No
radiative corrections are included (the tree-level result only); the same
reference \parencite{BahcallKamionkowskiSirlin1995} that fixes the
$(g_L,g_R)$ convention above also quantifies the size of the radiative
corrections neglected here.

The three terms in the bracket have a direct physical reading: $g_L^2$ is
the pure left-handed (helicity-conserving in the massless-electron limit)
contribution, constant in $T/E_\nu$; $g_R^2(1-T/E_\nu)^2$ is the
right-handed contribution, suppressed at large recoil $T\to E_\nu$ (a
right-handed electron cannot absorb the full neutrino energy while
respecting angular-momentum conservation in the massless limit); and the
cross term $-g_Lg_R\,m_eT/E_\nu^2$ is a genuine electron-mass effect --
it vanishes identically as $m_e\to0$ -- so it is never negligible at the
sub-MeV/MeV recoils probed by solar-neutrino scintillator detectors, where
$m_e/E_\nu = \mathcal{O}(1)$.

\subsection*{Kinematic endpoint}
For a massless incident neutrino scattering elastically off an electron at
rest, two-body relativistic kinematics fixes the maximum recoil, reached
when the neutrino backscatters ($\theta_\nu=\pi$):
\begin{equation}
  \keyresult{T_{\max}(E_\nu) = \frac{2E_\nu^2}{m_e + 2E_\nu}.}
\end{equation}
In the ultrarelativistic limit $E_\nu\gg m_e$ this saturates at
$T_{\max}\to E_\nu$ (the neutrino can in principle transfer its entire
energy); at solar-neutrino energies ($E_\nu\lesssim$ few MeV, comparable to
$m_e=0.511\ \mathrm{MeV}$) the endpoint is measurably below $E_\nu$, which
is why the recoil spectrum from a fixed incident line (e.g.\ the $^7$Be
monochromatic solar line) is a smooth edge rather than a sharp step at
$T=E_\nu$.
\end{theorybox}

\begin{implbox}
\begin{itemize}
  \item \pyfunc{NUE\_COUPLINGS}, \pyfunc{NUMUTAU\_COUPLINGS}: the two
    $(g_L,g_R)$ pairs above, built once at import time from
    \code{util.constant.SIN2\_THETA\_W}.
  \item \pyfunc{dsigma\_dT(E\_nu\_MeV, T\_MeV, g\_L, g\_R)} implements
    \labelcref{eq:nue-dsdt} directly, returning $\mathrm{cm^2/MeV}$; the
    prefactor's $(\hbar c)^2$ (\code{util.constant.HBARC\_MeV\_m}, converted
    to $\mathrm{MeV\,cm}$) is what turns $G_F^2 m_e$ -- natural units,
    $\mathrm{MeV^{-3}}$ -- into a cross section per unit energy in laboratory
    units. Enforces $0\leq T\leq T_{\max}(E_\nu)$ via
    \code{torch.where}, returning exactly zero outside that range rather
    than an unphysical negative or complex value.
  \item \pyfunc{nue\_cross\_section\_grid(E\_nu\_grid\_MeV, T\_grid\_MeV)},
    \pyfunc{numutau\_cross\_section\_grid(...)}: pre-evaluate
    \pyfunc{dsigma\_dT} on the outer-product grid
    \code{(E\_nu\_grid\_MeV, T\_grid\_MeV)}, shape \code{(n\_E, n\_T)},
    under \code{torch.no\_grad()} (the cross section is a fixed physical
    input, not a fit parameter) -- exactly the \code{cross\_section\_e}/
    \code{cross\_section\_x} arguments expected by
    \pyfunc{detector.common.event\_rate.true\_observable\_spectrum}
    (\cref{ch:detector-common}).
\end{itemize}
\end{implbox}

\begin{notebox}
\begin{itemize}
  \item \textbf{Tree level only.} \emph{Formulation:} \labelcref{eq:nue-dsdt}
    is the leading-order electroweak result; no QED or electroweak
    radiative corrections are applied. \emph{Justification and validity:}
    radiative corrections to this process are known to shift the effective
    couplings at the percent level \parencite{BahcallKamionkowskiSirlin1995};
    adequate for the illustrative/validation-level event-rate predictions
    this project currently targets, but a precision comparison against a
    real published rate measurement would need to fold them in.
  \item \textbf{Free target electron.} \emph{Formulation:} the target
    electron is treated as free and at rest, with no atomic binding energy
    or momentum. \emph{Justification and validity:} an excellent
    approximation whenever the momentum transfer (set by $T_{\max}$ above)
    is large compared to the atomic binding energy of the target's outer
    electrons -- true for essentially the whole recoil range relevant to
    MeV-scale solar-neutrino detection in a liquid scintillator, except very
    close to $T\to0$, where atomic binding effects (not modelled here) can
    in principle become relevant.
  \item \textbf{$\nu_\mu$/$\nu_\tau$ universality.} \emph{Formulation:}
    \pyfunc{NUMUTAU\_COUPLINGS} is used identically for $\nu_\mu$ and
    $\nu_\tau$. \emph{Justification and validity:} exact in the Standard
    Model at this order, since both interact with the target electron
    through the NC amplitude alone, which is flavour-blind by construction;
    this is precisely what licenses \texttt{detector.common}'s two-channel
    $\Phi_e/\Phi_x$ split (\cref{ch:detector-common}) instead of a full
    three-flavour cross-section grid.
\end{itemize}
\end{notebox}
